\article{$\GL_2(\R)$ 的轨道积分}{Orbital integrals for $\GL_2(\R)$}{D. Shelstad}{1}{107--110}
本文依照 [3]，简要介绍 $\GL_2(\R)$ 上紧支撑光滑（$C^\infty$，smooth）函数的轨道积分（orbital integrals）的刻画。类似的论证也适用于 $\GL_2(\C)$ [3]。我们先以适合这一刻画的形式回顾这些积分的一些熟知性质，再说明证明；[3] 给出了更优美的表述。\footnote{本工作得到美国国家科学基金会项目 MCS76-08218 的部分资助。原文 AMS（MOS）学科分类（1970）：主要 22E30。}

在 $G=\GL_2(\R)$ 上固定一个不变 $4$ 形式 $\omega_G$。若 $T$ 是 $G$ 的 Cartan 子群（Cartan subgroup），取
\[
 \omega_T=C_T\frac{d\gamma_1\,d\gamma_2}{\gamma_1\gamma_2},
\]
其中 $\gamma_1,\gamma_2$ 是 $\gamma$ 的特征值，按某个次序排列（该次序由 $T$ 的一次对角化规定）；常数 $C_T$ 如下：当 $T$ 分裂（split）时 $C_T=1$，否则 $C_T=i$。若 $f\in C_c^\infty(G)$，且 $\gamma\in T_{\mathrm{reg}}$（即 $G$ 中位于 $T$ 内的正则元素之集），定义
\[
 \Phi_f^T(\gamma)=\int_{G/T}f(g\gamma g^{-1})\frac{dg}{dt},
\]
这里 $dg,dt$ 分别为 $\omega_G,\omega_T$ 所定义的 Haar 测度（Haar measures）。于是 $\Phi_f^T$ 是 $T_{\mathrm{reg}}$ 上良定义的 $C^\infty$ 函数，在 Weyl 群作用下不变，并且在 $T$ 的某个相对紧集之外为零。令 $Z$ 为 $G$ 的标量矩阵群；于是 $Z=T-T_{\mathrm{reg}}$。$\Phi_f^T$ 在 $z\in Z$ 附近的行为可描述如下：存在 $z$ 在 $T$ 中的一个 Weyl 群不变邻域 $N_z$，以及定义在 $N_z$ 上、在 Weyl 群作用下不变的 $C^\infty$ 函数 $\Lambda_f^0(z,\,\cdot\,)$ 和 $\Lambda_f^1(z,\,\cdot\,)$，使得
\begin{equation}
 \Phi_f^T(\gamma)=\Lambda_f^0(z,\gamma)+\Lambda_f^1(z,\gamma)|D(\gamma)|^{-1/2}
 \tag{1}
\end{equation}
对 $\gamma\in N_z\cap T_{\mathrm{reg}}$ 成立。这里
\[
 D(\gamma)=\frac{(\gamma_1-\gamma_2)^2}{\gamma_1\gamma_2},
\]
其中 $\gamma_1,\gamma_2$ 仍为 $\gamma$ 的特征值。函数 $\Lambda_f^0(z,\,\cdot\,)$ 和 $\Lambda_f^1(z,\,\cdot\,)$ 依赖于 $T$，但记号中省略这一点。注意，等式 (1) 唯一确定了 $\Lambda_f^i(z,\,\cdot\,)$ 及其所有导数在 $Z\cap N_z$ 上的限制。因此，只要 $z\in N_{z'}$，便可令 $\Lambda_f^i(z)=\Lambda_f^i(z',z)$；导数也有类似定义。此外：
\begin{enumerate}
\item[(a)] 若 $T$ 分裂，则可取 $\Lambda_f^0(\,\cdot\,,\,\cdot\,)\equiv0$。
\item[(b)] 若 $X_T$ 表示 $\mathfrak{gl}_2(\C)$ 的泛包络代数（universal enveloping algebra）中心中算子 $X$ 在 Harish-Chandra 同构下的像，那么对于每个 $z\in Z$，$X_T\Lambda_f^1(z)$ 与 $T$ 无关。
\end{enumerate}
要确定 $\Lambda_f^i(z,\,\cdot\,)$ 的导数在 $Z\cap N_z$ 上的限制，只需对每个 $X$ 计算 $X_T\Lambda_f^i(z)$。我们不需要显式写出 $X_T\Lambda_f^1(z)$；该值可表示为 $Xf$ 在 $\bigl(\begin{smallmatrix}z&1\\0&z\end{smallmatrix}\bigr)$ 的共轭类上适当定义的积分。若 $T$ 不分裂，则 $X_T\Lambda_f^0(z)=c_GXf(z)$，其中常数 $c_G$ 仅依赖于 $G$ 上 Haar 测度的选择。

下面回顾证明。只需考虑对角群 $A$ 及如下 Cartan 子群：
\[
 B=\left\{\begin{pmatrix}a&b\\-b&a\end{pmatrix}:a^2+b^2\ne0\right\}.
\]
将 $B_{\mathrm{reg}}$ 中的元素写成
\[
 \gamma(\lambda,\theta)=\begin{pmatrix}\lambda\cos\theta&\lambda\sin\theta\\-\lambda\sin\theta&\lambda\cos\theta\end{pmatrix},
 \qquad\lambda>0,\quad\theta\not\equiv0\pmod\pi,
\]
这样写较为方便。经形式计算，可选择 $\omega_G$ 使得
\begin{equation}
 |D(\gamma)|^{1/2}\Phi_f^A(\gamma)
 =\frac12\left|\frac{\gamma_2}{\gamma_1}\right|^{1/2}\int_N f_0(n\gamma)\,dn,
 \tag{2}
\end{equation}
其中
\[
 \begin{gathered}
 N=\left\{\begin{pmatrix}1&x\\0&1\end{pmatrix}\right\},\qquad dn=dx,
 \qquad f_0(x)=\int_{K_0}f(kxk^{-1})\,dk,\\
 K_0=\left\{\begin{pmatrix}\cos\theta&\sin\theta\\-\sin\theta&\cos\theta\end{pmatrix}\right\},\qquad dk=d\theta;
 \end{gathered}
\]
并且
\begin{equation}
 \Phi_f^B(\gamma(\lambda,\theta))=
 \frac14\int_0^\infty
 \bigl(H_\lambda(e^t\theta,e^{-t}\theta)+H_\lambda(-e^t\theta,-e^{-t}\theta)\bigr)(e^t-e^{-t})\,dt,
 \tag{3}
\end{equation}
其中
\[
 H_\lambda(u,v)=\int_{K_0} f\left(\lambda k\exp\begin{pmatrix}0&u\\-v&0\end{pmatrix}k^{-1}\right)\,dk.
\]
式 (2) 来自 $\SL_2(\R)$ 的 Iwasawa 分解（Iwasawa decomposition），式 (3) 来自 Cartan 分解（Cartan decomposition）。沿用 [1] 的方法分析 (3) 右端的函数，可得 (1) 中的展开；余下的论证便可顺次完成。为计算 $X_T\Lambda_f^i(z)$，注意
\[
 |D|^{1/2}\Phi_{Xf}^T=X_T\bigl(|D|^{1/2}\Phi_f^T\bigr).
\]
这本质上就是 Harish-Chandra 的公式 $F_{Xf}^T=X_TF_f^T$ [2]；$|D|^{1/2}\Phi_f^A$ 即函数 $F_f^A$，而
\[
 (e^{i\theta}-e^{-i\theta})\Phi_f^B(\gamma(\lambda,\theta))=F_f^B(\gamma(\lambda,\theta)).
\]

下面证明这些性质足以刻画轨道积分。假定对每个 Cartan 子群 $T$，给定 $T_{\mathrm{reg}}$ 上定义的 $C^\infty$ 函数 $\Phi^T$，它在 Weyl 群作用下不变，并且在 $T$ 中某个相对紧集之外为零。又假定当 $T,T'$ 共轭时，$\Phi^T,\Phi^{T'}$ 满足显然应有的相容条件。最后，假定对每个 $T$ 及 $z\in Z$，存在 $z$ 在 $T$ 中的 Weyl 群不变邻域 $N_z$，以及其上的 Weyl 群不变 $C^\infty$ 函数 $\Lambda^0(z,\,\cdot\,)$、$\Lambda^1(z,\,\cdot\,)$，使得
\begin{equation}
 \Phi^T(\gamma)=\Lambda^0(z,\gamma)+\frac{\Lambda^1(z,\gamma)}{|D(\gamma)|^{1/2}},
 \qquad\gamma\in N_z\cap T_{\mathrm{reg}}.
 \tag{4}
\end{equation}
假定这些 $\Lambda^i(z,\,\cdot\,)$ 具有以下两条性质：
\begin{enumerate}
\item[(a)] 当 $T$ 分裂时，$\Lambda^0(\,\cdot\,,\,\cdot\,)\equiv0$；
\item[(b)] 对 $\mathfrak{gl}_2(\C)$ 泛包络代数中心中的每个 $X$，$X_T\Lambda^1(z,\,\cdot\,)$ 在 $Z\cap N_z$ 上的限制与 $T$ 无关。
\end{enumerate}
那么 [3] 的引理 4.1 断言：存在 $f\in C_c^\infty(G)$，使每个 $T$ 均满足 $\Phi_f^T=\Phi^T$。下面概述论证。

令 $G_r=G-Z$，$Y=\{(x_1,x_2)\in\R^2:x_2\ne0\}$。定义 $\pi:G_r\to Y$ 为 $\pi(g)=(\operatorname{trace}g,\det g)$；$\pi$ 是浸没（submersion），每条纤维都是 $G$ 中的一个共轭类。令 $S=\{(2x_1,x_1^2):x_1\ne0\}$。让 $T$ 遍历 $G$ 的 Cartan 子群，以下式定义 $Y-S$ 上的函数 $\psi$：
\[
 \psi(\pi(\gamma))=|D(\gamma)|^{1/2}\Phi^T(\gamma),\qquad\gamma\in T_{\mathrm{reg}}.
\]
若 $\psi$ 在 $S$ 附近为零，即每个 $\Phi^T$ 在 $Z$ 的某个邻域中为零，则容易找到 $f\in C_c^\infty(G)$ 使所有 $T$ 都满足 $\Phi_f^T=\Phi^T$（利用覆盖 $T_{\mathrm{reg}}\times G/T\to T_{\mathrm{reg}}^G$）。现假定 $\psi$ 能延拓为 $Y$ 上的光滑函数。由于 $\psi$ 具有紧支撑，可在 $Y$ 上使用单位分解（partition of unity），并假设 $\psi$ 的支撑包含于 $S$ 中某点的一个邻域（该邻域稍后指定）。

固定 $a\in S$ 及 $g\in\pi^{-1}(a)$。可在 $G_r$ 中找到 $g$ 的邻域 $N_1$，以及其上的坐标系 $y_1,\ldots,y_4$，使 $y_1=x_1\circ\pi$、$y_2=x_2\circ\pi$；不妨假定 $(y_i)$ 将 $N_1$ 映至 $\R^4$ 中的一个立方体。令 $N_2=\pi(N_1)$，并假定 $\psi\in C_c^\infty(N_2)$。将微分形式 $|x_2|^{-3/2}dx_1dx_2$ 拉回到 $N_1$。利用它与不变形式 $\omega_G$，在 $\pi$ 的每条纤维上构造一个 $G$ 不变测度。容易找到 $f\in C_c^\infty(N_1)$ 使
\[
 \int_{\pi^{-1}(x)}f=\psi(x),\qquad x\in N_2.
\]
另一方面，若 $\gamma\in N_1\cap T$ 且 $x=\pi(\gamma)$，则对 $h\in C_c^\infty(N_1)$ 有
\[
 \int_{\pi^{-1}(x)}h=|D(\gamma)|^{1/2}\Phi_h^T(\gamma).
\]
这是直接的坐标计算。因此 $\psi(x)=|D(\gamma)|^{1/2}\Phi_f^T(\gamma)$，$\psi$ 光滑时的论证至此完成。

接着注意，若与 Cartan 子群 $B$ 对应的函数 $\Lambda^0(z,\,\cdot\,)$ 对每个 $n$ 满足
\[
 \lim_{\theta\to0;\,(\theta\to\pi)}\frac{d^n}{d\theta^n}
 \bigl(|D(\gamma(\lambda,\theta))|^{1/2}\Lambda^0(\pm\lambda,\gamma(\lambda,\theta))\bigr)=0,
\]
或者更简单地说，对每个 $X,z',z$ 有 $X_B\Lambda^0(z',z)=0$，则 $\psi$ 可光滑地延拓到 $Y$。同前一样，省略 $z'$，只写 $X_B\Lambda^0(z)$。为处理一般的 $X_B\Lambda^0(z)$，我们转而使用速降函数及其轨道积分。

在 $G$ 或任何实约化群上，可引入 Harish-Chandra [2] 所定义的速降（rapidly decreasing，Schwartz）函数空间；[2] 连同其中列出的较早论文，对此空间上的“$F_f$”（正规化轨道积分）变换作了广泛分析。若 $f$ 在 $G$ 上速降，则 $\Phi_f^T$ 具有前面列举的性质，只是关于 $\Phi_f^T$ 支撑的陈述，须换成 $|D|^{1/2}\Phi_f^T$“在 $T_{\mathrm{reg}}$ 上速降”[2]。作刻画时，仍取上述 $\{\Phi^T\}$，但只要求 $|D|^{1/2}\Phi^T$ 在 $T_{\mathrm{reg}}$ 上速降。由于许多速降函数的轨道积分可以显式计算 [2]，论证思路较清楚，但细节较长。具体步骤如下。考虑 $A$ 上的速降函数 $|D|^{1/2}\Phi^A$。利用波包（wave-packets）[2] 找到 $f_1$，使 $F_{f_1}^A=|D|^{1/2}\Phi^A$；于是 $\Phi_{f_1}^A=\Phi^A$。再考虑\nobreak\footnote{\textbf{译者注（本次校订）：}这里须统一 $F_f^B$ 的正规化。前面已定义 $F_f^B=(e^{i\theta}-e^{-i\theta})\Phi_f^B=2i\sin\theta\,\Phi_f^B$，因此下式及随后定义 $F_{f_2}^B$ 的式子中，$\sin\theta\,\Phi^B$ 都应换成 $2i\sin\theta\,\Phi^B$。这样最后 $f=f_1+f_2$ 才满足 $\Phi_f^B=\Phi^B$。}
\[
 \sin\theta\,\Phi^B(\gamma(\lambda,\theta))-F_{f_1}^B(\gamma(\lambda,\theta)).
\]
由 (1) 可知，该函数在 $B$ 上为 $C^\infty$（并且作为 $\lambda$ 的函数速降）。然后主要利用 $\{x\in G:|\det x|=1\}$ 的离散系列表示（discrete series representations）的矩阵系数（matrix coefficients），找到 $f_2$，使
\[
 F_{f_2}^B=\sin\theta\,\Phi^B-F_{f_1}^B,\qquad F_{f_2}^A=0.
\]
令 $f=f_1+f_2$，则每个 $T$ 都满足 $\Phi_f^T=\Phi^T$。（对一般群，这个论证只刻画稳定轨道积分（stable orbital integrals）。）详情见 [4]。

回到原来的函数族 $\{\Phi^T\}$，其中 $\Phi^T$ 在 $T$ 中某个相对紧集之外为零。可找到速降函数 $f$ 使 $\Phi_f^T=\Phi^T$。于是对 $z\in Z$ 有 $X_B\Lambda^0(z)=c_GXf(z)$。必要时把 $f$ 乘以行列式的适当函数，可假设
\[
 \{x\in\R^\times:x=\det g,\ f(g)\ne0\}
\]
在 $\R^\times$ 中相对紧。于是可在 $C_c^\infty(G)$ 中找到函数 $f_1$，它在 $Z$ 的某个邻域上与 $f$ 一致。因而对所有 $z,X$ 均有 $Xf(z)=Xf_1(z)$，并且与 $\{\Phi^T-\Phi_{f_1}^T\}$ 对应的函数 $\psi$ 在 $Y$ 上光滑。再应用前面的论证，便完成了刻画定理的证明。

最后，固定 $Z$ 上的一个拟特征（quasi-character）$\chi$。假定 $f\in C^\infty(G)$ 满足 $f(zg)=\chi(z)f(g)$（$z\in Z,g\in G$），且支撑模 $Z$ 为紧。于是 $\Phi_f^T$ 良好定义，并具有前述性质，其中的支撑条件和 $Z$ 作用下的变换条件作相应调整即可。为刻画 $\{\Phi_f^T\}$，可如下论证。设 $\{\Phi^T\}$ 具有那些性质。容易找到 $T_{\mathrm{reg}}$ 上的 $C^\infty$ 函数 $\Phi_0^T$，它在 Weyl 群作用下不变，在 $T$ 的某个相对紧集之外为零，满足 (4)，并且
\[
 \Phi^T(\gamma)=\int_Z\chi(z^{-1})\Phi_0^T(z\gamma)\,dz,\qquad\gamma\in T_{\mathrm{reg}}.
\]
例如，可取 $\Phi_0^T$ 为 $\Phi^T$ 乘以 $|\det|$ 的适当函数。选取 $f_0\in C_c^\infty(G)$ 使每个 $T$ 都有 $\Phi_{f_0}^T=\Phi_0^T$。于是，由
\[
 f(g)=\int_Z\chi(z^{-1})f_0(zg)\,dz,\qquad g\in G
\]
定义的 $f$ 满足每个 $T$ 都有 $\Phi_f^T=\Phi^T$，且具有所要求的形式。

\begin{thebibliography}{4}
\bibitem[1]{c06-1} Harish-Chandra, \textit{Some results on an invariant integral on a semi-simple Lie algebra}, Ann. of Math. (2) \textbf{80} (1964), 457--508.
\bibitem[2]{c06-2} Harish-Chandra, \textit{Harmonic analysis on real reductive groups. I, II, III}, J. Functional Analysis \textbf{19} (1975), 104--204; Invent. Math. \textbf{36} (1976), 1--55; Ann. of Math. (2) \textbf{104} (1976), 117--201.
\bibitem[3]{c06-3} R. P. Langlands, \textit{Base change for GL(2)}（待刊）.
\bibitem[4]{c06-4} D. Shelstad, \textit{Characters and inner forms of a quasi-split group over R}, Compositio Math.（待刊）.
\end{thebibliography}
\noindent Columbia University（哥伦比亚大学）。
